\chapter{Introduction}

This chapter gives clear understanding on the introduction of our project like overview, motivation, objectives to achieve, methodology to implement and the outcome corresponding to objectives.

\section{Thesis/Project Overview}
Pharmacies rely on manual inventory and sales processes. Fragmented tracking causes stock mismanagement and financial waste. They are rely on human experience based forecasting for fututre sales prediction to control the rate of purchase of new medicines. Sales tracking and invoice generation are manual in most of the pharmacies in Bangladesh.


Our project, Pharmaciano is made to solve all the problems mentioned above. We implemented to such a system so that pharmacy can now automate the inventory and stock management. It provides AI based demand forecasting that helps to predict how the sales is going to be in the following months that is crusial to buy the stock for next. It is made by focusing on Bangladesh perspective, so it is easier for Bangladesh's pharmacy to convert existing manual system to ERP. We used Next.js for frontend, Express.js for backend, MongoDB for database, Redis for caching, Flutter for mobile app. Overall we made a robust solution for pharmacies in Bangladesh.


\section{Motivation}
In Bangladesh, most of the pharmacies are facing stock and expiry mismanagement as they make sales in manual process. As a result we the consumers sometime can not get medicines in time or the medicines which are available have expired. And many of the pharmacies do not provide us a detailed invoice that can give a customer about the transparency for the money they are spending. These issues drew our attention to solve the whole problem. This how we motivated to develop such an ERP system for pharmacy. 

\section{Objectives}
We have listed some goals for our project to be achieved. Here are these:

\begin{enumerate} 
    \item To develop a responsive and user-friendly UI
    \item To implement secure and scalable backend architecture
    \item To implement an inventory management system for medicines
    \item To maintain multiple branch functionality
    \item To manage multiple role access
    \item To generate automatic invoice
    \item To implement AI based predictions
\end{enumerate}
\section{Methodology}
We implemented the objectives of our project by using an "API-First" approach. For user-friendly and responsive UI we used Next.js and Shadcdn on frontend. We used Express.js on backend, MongoDB as database, flutter for mobile application. Frontend and Backend are communicating each other with RESTful API. To make the system secure we implemented JWT Token on protected API endpoints. 
\\
\begin{figure}[htbp]
\centering 
\begin{tikzpicture}[scale=2]
\node[draw, minimum height=29pt, align=center, minimum width=78pt] (node2) at (0.06,2.49) {Next.js\\based frontend};
\node[draw, minimum height=29pt, align=center] (node1) at (3.77,2.49) {Express.js\\based backend};
\node[draw, minimum height=29pt, align=center, minimum width=68pt] (node4) at (3.77,0.89) {MongoDB\\database};
\node[draw, minimum height=29pt, align=center, minimum width=77pt] (node3) at (0.04,0.89) {Flutter\\application};
  \draw[arrows=Latex-Latex] (node1.west) -- (node2.east);
  \draw[decorate, decoration=bent, arrows=Latex-Latex] (node1.south west) -- (node3.north);
  \draw[arrows=Latex-Latex] (node1.south) -- (node4.north);
  \node[draw=none, node font={\fontsize{9.5pt}{3.5pt}\selectfont}] at (1.93,2.6) {RESTful API};
  \node[draw=none, node font={\fontsize{9.5pt}{3.5pt}\selectfont}, rotate=30] at (2.78,1.73) {RESTful API};
  \node[draw=none, node font={\fontsize{9.5pt}{3.5pt}\selectfont}, rotate=19] at (1.37,1.73) {RESTful API};
  \node[draw=none, node font={\fontsize{9.5pt}{4.5pt}\selectfont}] at (4.20,1.68) {Mongoose};
  \node[draw, align=center, minimum width=79pt, minimum height=29pt] (node5) at (2,0.89) {AI Model};
  \draw[arrows=Latex-Latex] (node1.south) -- (node5.north);
  
\end{tikzpicture}
    \caption{Flow diagram}
    \label{fig:basicFlowDiagram}
\end{figure}
\\
This how we implemented an inventory management system to sale the medicines. We added multi-branch functionality by maintaining all the sales branch wise. We implemented Role Based Access Control(RBAC) for our REST APIs, this how we managed multiple role access like super admin, salesman, cashier e.t.c. We implemented automatic invoice generation by using \textit{jspdf} library on frontend. And we implemented AI based prediction by using OpenAI standard AI API as Google Gemini. 
\section{Project Outcome}
The outcomes achieved in this project are as follows:

\begin{enumerate}
    \item \textbf{Achieved a fully responsive user friendly UI} that works perfectly on desktop, laptop and mobile devices. That provides less effort to learn the system across all the stuff of branches.
    \item \textbf{Implemented a secure and scalable backend architecture} by utilizing secure authentication protocol, JWT and modular approach. That gives us flexibility to scaleup our system.
    \item \textbf{Successfully deployed full functional inventory management system for medicines} that provides visibility for proper stock levels.
    \item \textbf{Enabled multi-branch functionality} by implementing centralized data, saperation of each sales by branch.
    \item \textbf{Implemented Role Based Access Control(RBAC)} that manages permissions for admins, managers, cashier and sales Staff that provides strict data security and prevention of unauthorized access to sensitive operations.
    \item \textbf{Automated the entire billing and invoicing generation} that reduced billing turnaround time and eliminated manual arithmetic errors during checkout.
    \item \textbf{Integrated an AI-driven predictive analytics module}  that forecasts future medicine demand based on historical sales data, that significantly reducing instances of overstocking and understocking.
\end{enumerate}
\section{Organization of the Report}
In the following chapters we are going to discuss Background Study on second chapter, System Architecture on third chapter, Implementation and Results on fourth chapter, Standards and Design Constraints on fifth chapter and finally Conclusion on sixth chapter. 
